% !TeX root = ../main-german.tex
% Dokumentklasse, Pakete und zugehörige Einstellungen in Ladereihenfolge.

% Lokale Pakete im Unterordner tex/ finden.
\makeatletter
\def\input@path{{tex/}}
\makeatother

\RequirePackage{kvoptions-patch}
\documentclass[
  numbers=autoendperiod,
  ngerman,
  a4paper,
  twoside,
  bibliography=totoc,
  headsepline,
  cleardoublepage=empty,
  parskip=half,
  draft=false
]{scrbook}
\usepackage{scrlayer-scrpage}
\usepackage{iftex}

\usepackage{ifplatform}

% backticks (`) werden als solches in verbatim-Umgebungen dargestellt
% Details unter:
%   - https://tex.stackexchange.com/a/341057/9075
%   - https://tex.stackexchange.com/a/47451/9075
%   - https://tex.stackexchange.com/a/166791/9075
\usepackage{upquote}

% Setze Deutsch als Sprache
\usepackage[english,main=ngerman]{babel}
% Neue deutsche Trennmuster
\babelprovide[hyphenrules=ngerman-x-latest]{german}
%
% Hinweis von http://tex.stackexchange.com/a/321066/9075
% Ermögliche die Benutzung von "= als Trennstriche
\addto\extrasenglish{\languageshorthands{ngerman}\useshorthands{"}}

% Links verhalten sich so, wie sie sollen
% Zeilenumbrüche bei URLs auch bei Bindestrichen erlauben, auch wenn es verwirrend sein könnte: Gehört der Bindestrich zur URL oder ist es ein Trennstrich?
% Siehe https://tex.stackexchange.com/a/3034/9075.
\usepackage[hyphens]{url}
% \urlstyle{same}
%
% Hinweis von http://tex.stackexchange.com/a/10419/9075.
\makeatletter
\g@addto@macro{\UrlBreaks}{\UrlOrds}
\makeatother

%math stuff
\usepackage[
  centertags,    % (default) center tags vertically
  % tbtags,        % 'Top-or-bottom tags': For a split equation, place equation numbers level with the last (resp. first) line, if numbers are on the right (resp. left).
  sumlimits,    % (default) Place the subscripts and superscripts of summation symbols above and below
  % nosumlimits,   % Always place the subscripts and superscripts of summation-type symbols to the side, even in displayed equations.
  intlimits,     % Like sumlimits, but for integral symbols.
  % nointlimits,   % (default) Opposite of intlimits.
  namelimits,    % (default) Like sumlimits, but for certain 'operator names' such as det, inf, lim, max, min, that traditionally have subscripts placed underneath when they occur in a displayed equation.
  % nonamelimits,  % Opposite of namelimits.
  % leqno,         % Place equation numbers on the left.
  % reqno,         % Place equation numbers on the right.
  fleqn,         % Position equations at a fixed indent from the left margin rather than centered in the text column.
]{amsmath}
\SetMathAlphabet{\mathcal}{normal}{OMS}{amsa}{m}{n} %% AMS font for mathcal

% Provide \mathbb (blackboard bold, e.g. for number sets) when the active math
% font does not already define it. The thesis math fonts (mathdesign / newtxmath /
% unicode-math) supply it -- and amssymb clashes with mathdesign -- so only load
% amssymb as a fallback when \mathbb is still missing.
\makeatletter
\@ifundefined{mathbb}{\usepackage{amssymb}}{}
\makeatother

%%% Doc: http://mirror.ctan.org/tex-archive/macros/latex/contrib/mh/doc/mathtools.pdf
% Extends amsmath and fixes some bugs
\usepackage[fixamsmath,disallowspaces]{mathtools}

%%% Doc: http://www.ctan.org/info?id=fixmath
% LaTeX's default style of typesetting mathematics does not comply
% with the International Standards ISO31-0:1992 to ISO31-13:1992
% which indicate that uppercase Greek letters always be typeset
% upright, as opposed to italic (even though they usually
% represent variables) and allow for typesetting of variables in a
% boldface italic style (even though the required fonts are
% available). This package ensures that uppercase Greek be typeset
% in italic style, that upright $\Delta$ and $\Omega$ symbols are
% available through the commands \upDelta and \upOmega; and
% provides a new math alphabet \mathbold for boldface
% italic letters, including Greek.
\usepackage{fixmath}

%for theorems, replacement for amsthm
\usepackage[amsmath,hyperref]{ntheorem}
%%% Doc: http://mirror.ctan.org/tex-archive/macros/latex/contrib/onlyamsmath/onlyamsmath.dvi
% Warnt bei Benutzung von Befehlen die mit amsmath inkompatibel sind.

% Probably not needed.
% \usepackage[
% 	all,
% 	warning
% ]{onlyamsmath}

%% !!! If you change the font, be sure that words such as "workflow" can
%% !!! still be copied from the PDF. If this is not the case, you have
%% !!! to use glyphtounicode. See comment at cmap package.
%%
%% Background: "workflow" contains "fl" which is a ligature, which in turn
%%             is rendered as one character in the PDF and needs to be split
%%             whily copying.

\ifluatex
  \usepackage[no-math]{fontspec}
  \usepackage{unicode-math}

  % See https://tug.org/FontCatalogue/texgyretermes/ for more information
  \setmainfont{texgyretermes}[
    Extension = .otf,
    UprightFont = *-regular,
    BoldFont = *-bold,
    ItalicFont = *-italic,
    BoldItalicFont = *-bolditalic,
    Ligatures=TeX
  ]
  % See https://tug.org/FontCatalogue/texgyreheros/ for more information
  \setsansfont[Scale=.9]{TeX Gyre Heros Regular}
  % shapely l, upright quotes
  % Normal scaling is too large --> thus, we use ",Scale=.9"
\setmonofont[StylisticSet={1,3},Scale=.9]{Inconsolatazi4}

  % Enable proper ligatures
  % For more information see https://ctan.org/pkg/selnolig
  % language "english" or "ngerman" is passed to selnolig by the document class
  \usepackage{selnolig}

\else
  \RequirePackage{newtxtext}
  \RequirePackage{newtxmath}
  \RequirePackage[zerostyle=b,scaled=.9]{newtxtt}

  % Has to be loaded AFTER any font packages. See https://tex.stackexchange.com/a/2869/9075.
  \usepackage[T1]{fontenc}
\fi

% Even more symbols
%\usepackage{stmaryrd} % for \ovee, \owedge, \otimes
%\usepackage{marvosym} % for \Writinghand %patched to not redefine \Rightarrow
%\usepackage{mathrsfs} % nice cursive letters via \mathscr{}
%\usepackage{calrsfs} % slightly bolder \mathcal{} letters - does not look that good.

% Fallback font - if the subsequent font packages do not define a font (e.g., monospaced)
%     This is the modern package for "Computer Modern".
%     In case this gets activated, one has to switch from cmap package to glyphtounicode (in the case of pdflatex)
%\usepackage[%
%    rm={oldstyle=false,proportional=true},%
%    sf={oldstyle=false,proportional=true},%
%    tt={oldstyle=false,proportional=true,variable=true},%
%    qt=false%
%]{cfr-lm}

% Headings are typeset in Helvetica (which is similar to Arial)
% Font for the headings - overrides lmodern
%\usepackage[scaled=.95]{helvet}

% A calligraphic font would work, but is not required
%\usepackage{mathrsfs} %  \mathscr{ABC}

% Font for the main text
%     Font for the body text - overrides lmodern
%     Linux Libertine, see http://www.linuxlibertine.org/
%     Package parameter [osf] = lowercase (old-style) figures
%     rm = libertine in the body text, Linux Biolinum NOT as sans-serif font, but keep helvet (from above)
%\usepackage[rm]{libertine}

% Alternative Font: Palantino. It is recommended by Prof. Ludewig for German texts
%     Package parameter [osf] = lowercase (old-style) figures
%     Please only in German texts
%\usepackage{mathpazo} %ftp://ftp.dante.de/tex-archive/fonts/mathpazo/ - tip from DE-TEX-FAQ 8.2.1
% The euro sign
%     For Palatino (mathpazo.sty): proper euro sign
%     Alternative: \usepackage{eurosym}
% \newcommand{\EUR}{\ppleuro}

% Font for program code - overrides lmodern
%     If commented out, the default font lmodern is used
%     For typewriter-style keywords in the listings - does not work on old installations, since some font shapes (<>=) are missing
%\usepackage[scaled=.92]{luximono}
%\usepackage{courier}
% BeraMono as typewriter font, tip from http://tex.stackexchange.com/a/71346/9075
%\usepackage[scaled=0.83]{beramono}

\usepackage{setspace}
% Alternative package: https://ctan.org/pkg/leading

% Check and cross symbols
\usepackage{pifont}
%for scaling see http://tex.stackexchange.com/a/130236/9075

% Even more symbols
%\usepackage{stmaryrd} % for \ovee, \owedge, \otimes
%\usepackage{marvosym} % for \Writinghand %patched to not redefine \Rightarrow
%\usepackage{mathrsfs} % nice cursive letters via \mathscr{}
%\usepackage{calrsfs} %\mathcal{} ein bisserl dickeren buchstaben erzeugen - sieht net so gut aus.

\automark[section]{chapter}
\setkomafont{pageheadfoot}{\normalfont\sffamily}
\setkomafont{pagenumber}{\normalfont\sffamily}

\ihead[]{}
\chead[]{}
\ohead[]{\headmark}
\cfoot[]{}
\ofoot[\usekomafont{pagenumber}\thepage]{\usekomafont{pagenumber}\thepage}
\ifoot[]{}

% Optischer Randausgleich und Grauwertkorrektur. Siehe See http://www.ctan.org/tex-archive/macros/latex/contrib/microtype/

\usepackage[
  babel=true,
  expansion=alltext,
  protrusion=alltext-nott,
  final
]{microtype}

% \texttt{test -- test} - diese Einstellung behält "--" bei (und konveriert sie nicht zu einem Bindestrich)
\DisableLigatures{encoding = T1, family = tt* }

% tracking=true muss als Parameter des microtype-packages mitgegeben werden
% Deaktiviert, da dies bei Algorithmen seltsam aussieht

%\DeclareMicrotypeSet*[tracking]{my}{ font = */*/*/sc/* }%

% Hier wird festgelegt, dass alle Passagen in Kapitälchen automatisch leicht gesperrt werden.
% Quelle: http://homepage.ruhr-uni-bochum.de/Georg.Verweyen/pakete.html
% Deaktiviert, da sonst "BPEL", "BPMN" usw. wirklich komisch aussehen.
% Macht wohl nur bei geisteswissenschaftlichen Arbeiten Sinn.
%\SetTracking{ encoding = *, shape = sc }{ 45 }

% Avoid orphans and widows (a single line of a paragraph left at the bottom or top
% of a page). Using a package feels less hacky than tuning \clubpenalty and
% \widowpenalty manually.
\usepackage[all,defaultlines=3]{nowidow}

% Avoid a widow line being split off a displayed equation.
\displaywidowpenalty=10000

% Fuer deutsche Texte: Weniger Silbentrennung durch mehr frei Raum zwischen den Worten
\setlength{\emergencystretch}{3em}

\usepackage{graphicx}

% Base folder, so there is no need to repeat this over and over again.
\graphicspath{ {figures/} }

%%% Doc: http://mirror.ctan.org/tex-archive/macros/latex/contrib/pdfpages/pdfpages.pdf
\usepackage{pdfpages} % Include pages from external PDF documents in LaTeX documents

% Diagonal lines in a table - http://tex.stackexchange.com/questions/17745/diagonal-lines-in-table-cell
% Slashbox is not available in texlive (due to licensing) and also gives bad results. Thus, we use diagbox
\usepackage{diagbox}

\ifluatex
  \usepackage{spelling}
  \spellingoutput{off}
\fi

\usepackage[dvipsnames, table]{xcolor}

% Code Listings
\usepackage{listings}

% accsupp hält den Pfeil (siehe postbreak unten) aus dem kopierten Text heraus
\usepackage{accsupp}

\definecolor{eclipseStrings}{RGB}{42,0.0,255}
\definecolor{eclipseKeywords}{RGB}{127,0,85}
\colorlet{numb}{magenta!60!black}

% JSON definition
% Source: https://tex.stackexchange.com/a/433961/9075

\lstdefinelanguage{json}{
  basicstyle=\normalfont\ttfamily,
  commentstyle=\color{eclipseStrings}, % style of comment
  stringstyle=\color{eclipseKeywords}, % style of strings
  numbers=left,
  numberstyle=\scriptsize,
  stepnumber=1,
  numbersep=8pt,
  showstringspaces=false,
  breaklines=true,
  frame=lines,
  % backgroundcolor=\color{gray}, %only if you like
  string=[s]{"}{"},
  comment=[l]{:\ "},
  morecomment=[l]{:"},
  literate=
    *{0}{{{\color{numb}0}}}{1}
    {1}{{{\color{numb}1}}}{1}
    {2}{{{\color{numb}2}}}{1}
    {3}{{{\color{numb}3}}}{1}
    {4}{{{\color{numb}4}}}{1}
    {5}{{{\color{numb}5}}}{1}
    {6}{{{\color{numb}6}}}{1}
    {7}{{{\color{numb}7}}}{1}
    {8}{{{\color{numb}8}}}{1}
    {9}{{{\color{numb}9}}}{1}
}

\lstset{
  % everything between (* *) is a latex command
  escapeinside={(*}{*)},
  %
  language=json,
  %
  showstringspaces=false,
  %
  basicstyle=\footnotesize\ttfamily,
  %
  commentstyle=\slshape,
  %
  % Default: \rmfamily, damit werden die Strings im Quellcode hervorgehoben. Zusaetzlich evtl.: \scshape oder \rmfamily durch \ttfamily ersetzen. Dann sieht's aus, wie bei fancyvrb
  stringstyle=\ttfamily,
  %
  breaklines=true,            % Zeilen werden umbrochen
  %
  % Umbrochene Zeilen mit einem grauen Pfeil markieren. Dank accsupp wird der
  % Pfeil nicht in den kopierten Text übernommen. Quelle: https://tex.stackexchange.com/a/649475/9075
  postbreak=\raisebox{0ex}[0ex][0ex]{\BeginAccSupp{ActualText={}}\ensuremath{\color{gray}\hookrightarrow\space}\EndAccSupp{}},
  %
  breakatwhitespace=true,
  %
  % Alternative: fixed
  columns=flexible,
  %
  tabsize=2,                  % Groesse von Tabs
  %
  numbers=left,
  %
  numberstyle=\tiny,
  %
  basewidth=.5em,
  %
  xleftmargin=.5cm,
  %
  % aboveskip=0mm,
  %
  % belowskip=0mm,
  %
  captionpos=b
}

\ifpdftex
  % Ermögliche Umlaute falls \lstinputputlisting genutzt wird
  % Siehe https://stackoverflow.com/a/29260603/873282 und https://tex.stackexchange.com/a/24532/9075 für Details.
  % listingsutf8 hat im Juni 2020 nicht funktioniert.
  \lstset{extendedchars=true, literate=
      {á}{{\'a}}1 {é}{{\'e}}1 {í}{{\'i}}1 {ó}{{\'o}}1 {ú}{{\'u}}1
      {Á}{{\'A}}1 {É}{{\'E}}1 {Í}{{\'I}}1 {Ó}{{\'O}}1 {Ú}{{\'U}}1
      {à}{{\`a}}1 {è}{{\`e}}1 {ì}{{\`i}}1 {ò}{{\`o}}1 {ù}{{\`u}}1
      {À}{{\`A}}1 {È}{{\'E}}1 {Ì}{{\`I}}1 {Ò}{{\`O}}1 {Ù}{{\`U}}1
      {ä}{{\"a}}1 {ë}{{\"e}}1 {ï}{{\"i}}1 {ö}{{\"o}}1 {ü}{{\"u}}1
      {Ä}{{\"A}}1 {Ë}{{\"E}}1 {Ï}{{\"I}}1 {Ö}{{\"O}}1 {Ü}{{\"U}}1
      {â}{{\^a}}1 {ê}{{\^e}}1 {î}{{\^i}}1 {ô}{{\^o}}1 {û}{{\^u}}1
      {Â}{{\^A}}1 {Ê}{{\^E}}1 {Î}{{\^I}}1 {Ô}{{\^O}}1 {Û}{{\^U}}1
      {Ã}{{\~A}}1 {ã}{{\~a}}1 {Õ}{{\~O}}1 {õ}{{\~o}}1
      {œ}{{\oe}}1 {Œ}{{\OE}}1 {æ}{{\ae}}1 {Æ}{{\AE}}1 {ß}{{\ss}}1
      {ű}{{\H{u}}}1 {Ű}{{\H{U}}}1 {ő}{{\H{o}}}1 {Ő}{{\H{O}}}1
      {ç}{{\c c}}1 {Ç}{{\c C}}1 {ø}{{\o}}1 {å}{{\r a}}1 {Å}{{\r A}}1
  }
\fi

\lstloadlanguages{% Check dokumentation for further languages...
  %[Visual]Basic
  %Pascal
  %C
  %C++
  %XML
  %HTML
}

% For easy quotations: \enquote{text}
% This package is very smart when nesting is applied, otherwise textcmds (see below) provides a shorter command
\usepackage[autostyle=true]{csquotes}

% Enable using "`quote"' - see https://tex.stackexchange.com/a/150954/9075
\defineshorthand{"`}{\openautoquote}
\defineshorthand{"'}{\closeautoquote}

% bessere Abstaende innerhalb der Tabelle (Layout))
% -------------------------------------------------
% \toprule, \midrule, \bottomrule
% Doc: https://texdoc.org/serve/booktabs/0
\usepackage{booktabs}

% Extended enumerate, such as \begin{compactenum}
\usepackage{paralist}
\usepackage[
  backend       = biber, %biber does not work with 64x versions alternative: bibtex8; minalphanames only works with biber backend
  sortcites     = true,
  bibstyle      = alphabetic,
  citestyle     = alphabetic,
  giveninits    = true,
  useprefix     = false, %"von, van, etc." will be printed, too. See below.
  minnames      = 1,
  minalphanames = 3,
  maxalphanames = 4,
  maxbibnames   = 99,
  maxcitenames  = 2,
  natbib        = true,
  eprint        = true,
  url           = true,
  doi           = true, %source: http://tex.stackexchange.com/a/23118/9075
  isbn          = true, %source: http://tex.stackexchange.com/a/23118/9075
  backref       = true]{biblatex}

% enable more breaks at URLs. See https://tex.stackexchange.com/a/134281.
\setcounter{biburllcpenalty}{7000}
\setcounter{biburlucpenalty}{8000}

\addbibresource{bibliography/bibliography.bib}
%\addbibresource[datatype=bibtex]{\bibliography{bibliography}}

% Embed the bibliography data into the PDF so that readers can extract it
% (e.g., with JabRef via "Import into library" on the PDF)
\usepackage{embedfile}
\embedfile[mimetype=application/x-bibtex,desc=BibTeX entries of all references]{bibliography/bibliography.bib}

% Farbige Tabellen
% ----------------
% Das Paket colortbl wird inzwischen automatisch durch xcolor geladen
%
% Erweiterte Funktionen innerhalb von Tabellen
% --------------------------------------------
%%% Doc: http://mirror.ctan.org/tex-archive/macros/latex/contrib/multirow/multirow.sty
\usepackage{multirow} % Mehrfachspalten
%
%%% Doc: Documentation inside dtx Package
\usepackage{dcolumn}  % Ausrichtung an Komma oder Punkt

% Etwas mehr Luft zwischen den Tabellenzeilen (\extrarowheight wird von array bereitgestellt, das dcolumn lädt)
\setlength{\extrarowheight}{1pt}

%%% Doc: http://mirror.ctan.org/tex-archive/macros/latex/contrib/supertabular/supertabular.pdf
%\usepackage{supertabular}

% tikz / pgfplots stack used by the "Figures with tikz" and "Plots with pgfplots"
% examples. Loading pgfplots pulls in tikz.
\usepackage{pgfplots}
\pgfplotsset{compat=1.18}

% longtable: tables that automatically break across page boundaries.
\usepackage{longtable}

%%% Footnotes/Endnotes ===================================================

%%% Doc: http://mirror.ctan.org/tex-archive/macros/latex/contrib/footmisc/footmisc.pdf
%
\usepackage[
  bottom,      % Footnotes appear always on bottom. This is necessary specially when floats are used
  stable,      % Make footnotes stable in section titles
  % perpage,     % Reset on each page
  % para,        % Place footnotes side by side of in one paragraph.
  % side,        % Place footnotes in the margin
  ragged,      % Use RaggedRight
  % norule,      % Suppress rule above footnotes
  multiple,    % Rearrange multiple footnotes intelligent in the text.
  % symbol,      % Use symbols instead of numbers
]{footmisc}

\counterwithout{footnote}{chapter} % Continuous numbering of footnotes across chapters

\interfootnotelinepenalty=10000 % Prevents footnotes from continuing onto the facing page

% Put footnotes below floats
% Source: https://tex.stackexchange.com/a/32993/9075
\usepackage{stfloats}
\fnbelowfloat

% Extended support for footnotes

%\usepackage{dblfnote}  % Two-column footnotes
%
% No superscript numbers in the footnote (KOMA-Script-specific):
%\deffootnote[1.5em]{0pt}{1em}{\makebox[1.5em][l]{\bfseries\thefootnotemark}}
%
% Increase the spacing between footnotes:
%\setlength{\footnotesep}{.85\baselineskip}
%
% Following command disables the separating line of the footnote
%\renewcommand{\footnoterule}{}
%
%\addtolength{\skip\footins}{\baselineskip} % Spacing text <-> footnote

% Always place footnotes at the very bottom on a \raggedbottom page
% fnpos comes from the yafoot package
%\usepackage{fnpos}
%\makeFNbelow
%\makeFNbottom

% TODO (and comment) configuration
%
% - \todo (from todo, easy-todo, todonotes) / \TODO (from fixmetodonotes) - for "normal" TODOs
% - \todofix - "important" TODOs
%
% - \textcomment - highlights text and has a hover comment
% - \sidecomment - just puts a comment to the side. Note: \comment MUST NOT be used as command name, it is already defined by much packages (mathdesign, mindflow, verbatim, and others)
%
% - \missingfigure
%
% - \textmarker
% - \modified
% - \change      - adresses a review comment

% Enable nice comments
\usepackage{pdfcomment}

\usepackage[group-minimum-digits=4,per-mode=fraction]{siunitx}
\addto\extrasgerman{\sisetup{locale = DE}}

% See http://tex.stackexchange.com/a/83051/9075
% Normally, doesn't work with hyperref, but cleveref (or zref-vario) fixes that
\usepackage[ngerman]{varioref}


% Enable that parameters of \zcref{}, \ref{}, \cite{}, ... are linked so that a reader can click on the number an jump to the target in the document
\usepackage{hyperref}
\usepackage{graphicx}
% Enable hyperref without colors and without bookmarks
\hypersetup{
  hidelinks,
  colorlinks=true,       % Links erhalten Farben statt Kaeten
  raiselinks=true,       % calculate real height of the link
  allcolors=black,
  pdfstartview=Fit,
  breaklinks=true,       % Links ueberstehen Zeilenumbruch
  hypertexnames=false,   % Fix jumping to algorithm line - http://tex.stackexchange.com/a/156404/9075
}

% Enable correct jumping to figures when referencing
\usepackage[all]{hypcap}

% Hint by https://tex.stackexchange.com/a/193370/9075 to suppress strange outputs of the babel package
% Example strange output: Package babel Info: Redefining ngerman shorthand "|
\usepackage{etoolbox}
\makeatletter
\patchcmd{\@decl@short}{\bbl@info}{\@gobble}{}{}
\makeatother


%%%
% Allows rotating figures by 90 degrees
% Alternative package: rotating. However, that only rotates the image, whereas lscape also rotates the PDF page.
% The lscape package does not rotate the page either
\usepackage{pdflscape}

% Align floats to the top on float-only pages instead of vertically centering them.
%     See http://tex.stackexchange.com/a/28565/9075
\makeatletter
\setlength{\@fptop}{0pt}
\setlength{\@fpbot}{0pt plus 1fil}
\makeatother

\usepackage[caption=false,font=footnotesize]{subfig}

% Alternative for making subfigures:
% Part of the caption package. See http://www.ctan.org/pkg/caption
% Replaces the subfigure and subfig packages - see https://tex.stackexchange.com/a/13778/9075
%
% (subfigure is outdated. subfig is maintained, but subcaption is better)
% See: http://tex.stackexchange.com/questions/13625/subcaption-vs-subfig-best-package-for-referencing-a-subfigure
%\usepackage[hypcap=true]{subcaption}

\usepackage{mindflow}

% https://ctan.org/pkg/algorithms
% Consists of two environments: algorithm and algorithmic
% Although oudated, it defines the "algorithm" float enviornment
% TODO: Define floating environment "algorithm" in other ways
\usepackage[chapter]{algorithm}

% https://ctan.org/pkg/algpseudocodex
% Successor of algorithmicx; more modern than https://ctan.org/pkg/algorithms
\usepackage{algpseudocodex}

% Extensions for references inside the document (\zcref{fig:sample}, ...)
% Use \zcref{...} (and \Zcref{...} at the beginning of a sentence) instead of \ref: Type of reference included in the link
% That means, "Figure 5" is a full link instead of just "5".
% The names follow the document language (babel); zref-clever is compatible with tagged PDFs.
% Each \label also sets a zref label, so \zcref works with standard \label.
\usepackage{zref-clever}
\zcsetup{cap, nameinlink=true}
% Page-aware references: \zvref{...} (and \zvref[S]{...}) reads like "Figure 5 on the next page"
\usepackage{zref-vario}
\zcLanguageSetup{english}{type=page, Name-sg=page, Name-pl=pages} % "on page 5", not "on Page 5"

\usepackage{lipsum}

% For demonstration purposes only
% These packages can be removed when all examples have been deleted
\usepackage[math]{blindtext}
\usepackage{mwe}
\usepackage[realmainfile]{currfile}
\usepackage{tcolorbox}
\tcbuselibrary{listings}

\usepackage[
  translate=babel,
  abbreviations,         % create "abbreviations" glossary
  nomain,                % don't create "main" glossary
  stylemods=longbooktabs % do the adjustments for the longbooktabs styles
]{glossaries-extra}
\setglossarystyle{long3col-booktabs}

% Hint by https://tex.stackexchange.com/a/463188/9075
% \usepackage{glossary-longextra}

% Following is required if the abbreviation list should be sorted automatically (\printglossary / \printglossaries)
% Not required, if we printed the entries in-order (using \printunsrtglossaries)
% Required to have the German chapter name % Source: https://tex.stackexchange.com/a/426392/9075
\makeglossaries

% !TeX root = ../main-german.tex
% Note the difference between "acronym" and "abbreviation": https://english.stackexchange.com/q/3755/66058
% We regard abbreviation as superset of acronym
% `er` is used by the English example, `fr` by the German one. Both are defined in
% both abbreviation lists because the template ships a single abbreviations.tex that
% every main file (English and German) loads -- an undefined entry breaks the build.
\newabbreviation{er}{ER}{error rate}
\newabbreviation{fr}{FR}{Fehlerrate}
\newabbreviation[shortplural={RDBMS},longplural={RDBMS}]{rdbms}{RDBMS}{Relational Database Management System}
\newabbreviation[category=abbreviationdot]{dante}{DANTE e.V.}{Deutschsprachige Anwendervereinigung \TeX\ e.V.}
\newabbreviation{gp}{G.P.}{General Practitioner}

% Source: https://tex.stackexchange.com/a/154566/9075
\newabbreviation{html}{HTML}{hypertext markup language}
\newacronym{nato}{NATO}{North Atlantic Treaty Organization}



\usepackage{xspace}
% Macht \xspace und \enquote kompatibel
\makeatletter
\xspaceaddexceptions{\grqq \grq \csq@qclose@i \} }
\makeatother

\ifluatex
  % Enable correct rendering of ligatures - provided by https://ctan.org/pkg/autotype
  % See ADR-0008 for alternatives
  \usepackage{autotype}
\fi

% correct bad hyphenation here
\hyphenation{
  Spe-zi-fi-ka-tion
  In-te-gra-tion
  An-for-de-rung An-for-de-run-gen
  Be-nut-zer-ober-flä-che
  Mes-sung-en
  aus-zu-tau-schen
  Lauf-zeit-in-for-ma-tionen
  % May not be hypphenated
  AROMA TOSCA BPMN OASIS OMG DMTF IT DevOps
}

% lccaps provides \textlcc (lowercase caps), used by \initialism in commands.tex.
% Doc: https://ctan.org/pkg/lccaps
\usepackage{lccaps}


\usepackage[
  title={Shor's Algorithmus auf elliptischen und hyperelliptischen Kurven},
  author={Philipp Eckstein},
  type=master,
  institute=sec,
  course={Lehramt Informatik},
  examiner={Prof. Dr. Ralf Küsters},
  supervisor={Dr.\ Pascal Reisert},
  startdate={28.\ Juli 2026},
  enddate={13.\ Januar 2027}
]{scientific-thesis-cover}


\ifpdftex
  % Enable copy and paste of text from the PDF
  % Only required for pdflatex. It "just works" in the case of lualatex.
  % Alternative: cmap or mmap package
  % mmap enables mathematical symbols but does not work with the newtx font set
  % See: https://tex.stackexchange.com/a/64457/9075
  % Other solutions outlined at http://goemonx.blogspot.de/2012/01/pdflatex-ligaturen-und-copynpaste.html and http://tex.stackexchange.com/questions/4397/make-ligatures-in-linux-libertine-copyable-and-searchable
  % Troubleshooting outlined at https://tex.stackexchange.com/a/100618/9075
  %
  % According to https://tex.stackexchange.com/q/451235/9075 this is the way to go
  \input{glyphtounicode}
  \pdfgentounicode=1
\fi
% DM: line-breaking-description env vom daniel w.

% credit goes to daniel w. :-)
%% --- Descriptions with line breaks in labels ---------------------------------
\usepackage{calc}
\usepackage{xifthen}
\usepackage[inline]{enumitem}

% fix incompatibilities between KOMA and other packages, mainly float.
% should be loaded at the very end - see http://tex.stackexchange.com/a/156256/9075
\usepackage{scrhack}
